\answer{ex:19:1}{(a)~一个突触给出 $6.7\mV$（$R=66.7$~M$\Omega$），因此要能到达阈值至少需要 $4$ 个（三个只能渐近地达到）；$10\ms$ 内需要 $8$ 个；$5\ms$ 内需要 $14$ 个。(b)~$0.1\ms\ll\tau_m$，泄漏带来的修正 $\sim0.25\%$。}
\answer{ex:19:2}{(a)~$1.75\cdot10^{10}$（四分之一的混合器对什么都响应）、$4.4\cdot10^9$ 和 $0.06$（$k=40$ 时几乎从不响应）。(b)~$10^3$ 倍：$2\cdot10^5$ 对 $200$。}
\answer{ex:19:3}{(a)~由二项式系数的对称性，$C(2n,n)=2^{2n-1}$。(b)~$0.93$ 和 $0.09$；过渡区在 $P$ 上的宽度 $\sim\sqrt{2n}$，相对宽度 $\sim1/\sqrt n$。}
\answer{ex:19:4}{只靠一个树突，无论输入多少，都有 $V_{\text{胞体}}<\kappa E_E<\theta$；$g_0=0.5g_L$ 时每个树突上需要 $n\ge3$ 个输入。}
\answer{ex:19:5}{$I\ge C\sigma_V/\sigma_t=15$~nA，即 $150$ 个同步突触——不现实；小神经元只需 $1$~nA，而用一个 $4$~nA的突触，抖动为 $5$~\textmu s。}
\answer{ex:19:6}{近端的最大值：$0.03\,\tau_m$ 和 $0.97$（$L=0.2$）；$0.32\,\tau_m$ 和 $0.66$（$L=1$）；$0.79\,\tau_m$ 和 $0.33$（$L=2$）。按总电容所作的估计~\eqref{eq:dendriteload}~只在 $L\ll1$ 时成立。}
\answer{ex:19:7}{开放问题。$m$ 固定时，响应时间随记忆量 $P$ 线性增长；比例 $m=\varphi n$ 固定时，它不再依赖于 $n$，大神经元的缓慢是对许多输入求和的结果，而不是尺寸本身造成的。}
